\section{Parent Hamiltonian statements for the examples}\label{sec:example_parent_hamiltonians}
%% ===================================================================

The parent-Hamiltonian construction of \cite[lines~1988--2000]{Cirac2021Matrix},
reviewed in Chapter~\ref{ch:parent_hamiltonian}, is recorded here in the form
\begin{align}
    \mathcal G_L(A)
     & =
    \left\{\sum_{i_1,\ldots,i_L}
    \tr(A^{i_1}\cdots A^{i_L}X)
    \ket{i_1\cdots i_L}
    : X\in \MN{D}\right\}.
    \notag
\end{align}
Thus
\begin{align}
    \ker h_L  & = \mathcal G_L(A), \notag \\
    H_N^{(L)} & = \sum_i (h_L)_i.
    \notag
\end{align}
We write $\GS(H)=\ker(H-E_0(H)\Id)$ for the ground eigenspace.

\begin{theorem}[GHZ two-site parent space]\label{thm:ghz_ground_space_two}
    \lean{MPSTensor.ghz_groundSpace_two_eq_span}
    \leanok
    \uses{def:ghz_tensor, def:ground_space_parent, def:two_site_ket}
    For the GHZ tensor $A$ of Definition~\ref{def:ghz_tensor},
    \begin{align}
        \mathcal G_2(A)
         & =
        \spn_{\C}\{\ket{00},\ket{11}\}.
        \notag
    \end{align}
\end{theorem}

\begin{proof}
    \leanok
    \uses{def:ghz_tensor, def:ground_space_map, def:ground_space_parent, def:two_site_ket}
    By definition,
    $\Gamma_2(X)(\sigma_1,\sigma_2)
        =\tr(A^{\sigma_1}A^{\sigma_2}X)$.
    For the GHZ matrices $A^0=\diag(1,0)$ and $A^1=\diag(0,1)$, whenever
    $\sigma_1\ne\sigma_2$, one has $A^{\sigma_1}A^{\sigma_2}=0$. Hence
    $\Gamma_2(X)=X_{00}\ket{00}+X_{11}\ket{11}$, and the image of~$\Gamma_2$
    is $\spn_{\C}\{\ket{00},\ket{11}\}$.
\end{proof}

\begin{theorem}[Constant GHZ product vectors satisfy the cyclic constraints]
    \label{thm:constant_ket_mem_ghz_chain_ground_space}
    \lean{MPSTensor.constantKet_mem_ghz_chainGroundSpace}
    \leanok
    \uses{def:ghz_tensor, def:chain_ground_space, def:constant_ket, thm:ghz_ground_space_two}
    For $N\ge2$ and $a\in\{0,1\}$, one has
    $\ket{a}^{\otimes N}\in\mathcal G_{N,2}(A_{\mathrm{GHZ}})$.
\end{theorem}

\begin{proof}
    \leanok
    \uses{thm:ghz_ground_space_two}
    For a cyclic nearest-neighbour window and outside assignment $\tau$, set
    $R_{i,\tau}:=\left.\ket{a}^{\otimes N}\right|_{[i,i+1],\tau}$.
    For $b,c\in\{0,1\}$ with $b\ne c$, the completed configuration is not
    constant, hence $R_{i,\tau}(b,c)=0$. Therefore
    \begin{align}
        R_{i,\tau}
         & \in
        \{v:v(0,1)=v(1,0)=0\}
        =
        \spn_{\C}\{\ket{00},\ket{11}\}
        =
        \mathcal G_2(A_{\mathrm{GHZ}}),
        \notag
    \end{align}
    so $\ket{a}^{\otimes N}\in\mathcal G_{N,2}(A_{\mathrm{GHZ}})$.
\end{proof}

\begin{theorem}[GHZ periodic nearest-neighbour ground space]
    \label{thm:ghz_chain_ground_space_two}
    \lean{MPSTensor.ghz_chainGroundSpace_two_eq_constant_span}
    \leanok
    \uses{def:ghz_tensor, def:chain_ground_space, def:constant_ket, thm:ghz_ground_space_two}
    For $N\ge2$, the common nearest-neighbour cyclic ground space of the GHZ
    tensor satisfies
    \begin{align}
        \mathcal G_{N,2}(A_{\mathrm{GHZ}})
         & =
        \spn_{\C}\{
        \ket{0}^{\otimes N},
        \ket{1}^{\otimes N}
        \}.
        \notag
    \end{align}
    This is the cyclic-window form of the two-fold degeneracy in
    \cite[line~2205]{Cirac2021Matrix}.
\end{theorem}

\begin{proof}
    \leanok
    \uses{thm:ghz_ground_space_two, thm:constant_ket_mem_ghz_chain_ground_space}
    Let $\psi\in\mathcal G_{N,2}(A_{\mathrm{GHZ}})$. For every cyclic
    nearest-neighbour window and outside assignment $\tau$, one has
    \begin{align}
        \left.\psi\right|_{[i,i+1],\tau}
         & \in
        \mathcal G_2(A_{\mathrm{GHZ}})
        =
        \spn_{\C}\{\ket{00},\ket{11}\}.
        \notag
    \end{align}
    Taking the outside assignment from the word $\sigma$ gives
    $\left.\psi\right|_{[i,i+1],\sigma}
        (\sigma_i,\sigma_{i+1})=\psi(\sigma)$.
    Hence a mixed nearest-neighbour word has zero coefficient:
    $\sigma_i\ne\sigma_{i+1}\Longrightarrow\psi(\sigma)=0$.
    A nonconstant binary cyclic word has a mixed edge:
    $\sigma\notin\{(0,\ldots,0),(1,\ldots,1)\}
        \Longrightarrow\exists i,\ \sigma_i\ne\sigma_{i+1}$.
    Therefore
    \begin{align}
        \psi
         & =
        \psi(0,\ldots,0)\ket{0}^{\otimes N}
        +
        \psi(1,\ldots,1)\ket{1}^{\otimes N}.
        \notag
    \end{align}
    Conversely, Theorem~\ref{thm:constant_ket_mem_ghz_chain_ground_space}
    gives
    $\ket{0}^{\otimes N},\ket{1}^{\otimes N}
        \in\mathcal G_{N,2}(A_{\mathrm{GHZ}})$, and hence their span lies in
    $\mathcal G_{N,2}(A_{\mathrm{GHZ}})$.
\end{proof}

\begin{remark}[GHZ parent interaction]\label{rem:ghz_parent_interaction}
    For the GHZ tensor of Definition~\ref{def:ghz_tensor},
    \begin{align}
        \mathcal G_2(A)
         & = \spn\{\ket{00},\ket{11}\}, \notag \\
        h_{\mathrm{GHZ}}
         & = \ketbra{01}{01}+\ketbra{10}{10}
        = \Id-\ketbra{00}{00}-\ketbra{11}{11}.
        \notag
    \end{align}
    The periodic nearest-neighbour parent Hamiltonian therefore satisfies
    \begin{align}
        H_{\mathrm{GHZ},N}
         & = \sum_i (h_{\mathrm{GHZ}})_i, \notag              \\
        \ker H_{\mathrm{GHZ},N}
         & = \spn\{\ket{0}^{\otimes N},\ket{1}^{\otimes N}\},
        \notag
    \end{align}
    using the cyclic-window equality
    $\mathcal G_{N,2}(A_{\mathrm{GHZ}})
        =\spn\{\ket{0}^{\otimes N},\ket{1}^{\otimes N}\}$ from
    Theorem~\ref{thm:ghz_chain_ground_space_two}.
    This is the ferromagnetic Ising parent Hamiltonian
    \cite[lines~1194--1196 and line~2205]{Cirac2021Matrix};
    see also \cite[Section~3]{FernandezGonzalez2015Uncle} for the displayed
    two-site space and interaction.
\end{remark}

\begin{definition}[Cluster-state stabilizer]\label{def:cluster_stabilizer}
    \lean{MPSTensor.clusterStabilizer}
    \leanok
    The three-site stabilizer of \cite[local TeX lines~374--387]{PerezGarcia2007Matrix}
    is $K = \sigma^z_1\sigma^x_2\sigma^z_3$. In the computational basis, with
    $\sigma^z\ket{s} = (-1)^s\ket{s}$ and $\sigma^x\ket{s} = \ket{s+1}$, it acts
    on a three-site vector $v$ by
    \begin{align}
        (Kv)(s_1,s_2,s_3)
         & = (-1)^{s_1+s_3}\, v(s_1, s_2+1, s_3).
        \notag
    \end{align}
\end{definition}

\begin{theorem}[Cluster three-site parent space is the stabilizer eigenspace]
    \label{thm:cluster_ground_space_three_eigenspace}
    \lean{MPSTensor.cluster_groundSpace_three_eq_eigenspace}
    \leanok
    \uses{def:cluster_tensor, def:ground_space_parent, def:cluster_stabilizer}
    For the cluster tensor $A$ of Definition~\ref{def:cluster_tensor},
    \begin{align}
        \mathcal G_3(A)
         & = \{v : Kv = v\},
        \notag
    \end{align}
    the $+1$ eigenspace of $K = \sigma^z_1\sigma^x_2\sigma^z_3$.
\end{theorem}

\begin{proof}
    \leanok
    \uses{def:cluster_tensor, def:ground_space_map, def:cluster_stabilizer}
    Write $q = 1/\sqrt2$. The matrices $A^s = \ketbra{\pm}{s}$ have entries
    $A^s_{ij} = q(-1)^{is}\delta_{sj}$, so
    \begin{align}
        \Gamma_3(X)(s_1,s_2,s_3)
         & = \tr(A^{s_1}A^{s_2}A^{s_3}X)
        = q^3(-1)^{s_1s_2+s_2s_3}
        \bigl(X_{s_30} + (-1)^{s_1}X_{s_31}\bigr).
        \label{eq:cluster_gamma_three}
    \end{align}
    The sign $(-1)^{s_1s_2+s_2s_3}$ changes by $(-1)^{s_1+s_3}$ when $s_2$ is
    replaced by $s_2+1$, so $K\Gamma_3(X) = \Gamma_3(X)$ by
    (\ref{eq:cluster_gamma_three}). Conversely, if $Kv = v$, then
    $v(s_1,1,s_3) = (-1)^{s_1+s_3}v(s_1,0,s_3)$, and the matrix $X$ with rows
    $X_{s0} = v(0,0,s) + v(1,0,s)$, $X_{s1} = v(0,0,s) - v(1,0,s)$ satisfies
    $\Gamma_3(X) = 2q^3\,v$ by (\ref{eq:cluster_gamma_three}); hence
    $v = \Gamma_3(X/(2q^3)) \in \mathcal G_3(A)$.
\end{proof}

\begin{theorem}[The stabilizer has eigenvalue $-1$ on the source's matrices]
    \label{thm:cluster_source_ground_space_three_eigenspace}
    \lean{MPSTensor.clusterSourceTensor, MPSTensor.clusterSource_groundSpace_three_eq_eigenspace_neg_one}
    \leanok
    \uses{def:ground_space_parent, def:cluster_stabilizer}
    Let $A_{\mathrm{src}}$ be the tensor with the two matrices displayed in
    \cite[local TeX lines~378--387]{PerezGarcia2007Matrix}, the source's
    physical labels $1,2$ written as $0,1$:
    \begin{align}
        A_{\mathrm{src}}^0
         & = \begin{pmatrix} 0 & 0 \\ 1 & 1 \end{pmatrix},
        \qquad
        A_{\mathrm{src}}^1
        = \begin{pmatrix} 1 & -1 \\ 0 & 0 \end{pmatrix}.
        \notag
    \end{align}
    Then
    \begin{align}
        \mathcal G_3(A_{\mathrm{src}})
         & = \{v : Kv = -v\},
        \notag
    \end{align}
    the $-1$ eigenspace of $K = \sigma^z_1\sigma^x_2\sigma^z_3$.
\end{theorem}

\begin{proof}
    \leanok
    \uses{def:ground_space_map, def:cluster_stabilizer}
    With $A_{\mathrm{src}}^t = \ket{1-t}\bra{r_t}$, $r_0 = (1,1)$,
    $r_1 = (1,-1)$, one has $\braket{r_t}{u} = (-1)^{tu}$ and
    \begin{align}
        \Gamma_3(X)(t_1,t_2,t_3)
         & = (-1)^{t_1(1-t_2)+t_2(1-t_3)}
        \bigl(X_{0,1-t_1} + (-1)^{t_3}X_{1,1-t_1}\bigr).
        \label{eq:cluster_source_gamma_three}
    \end{align}
    Replacing $t_2$ by $t_2+1$ changes the exponent by $1+t_1+t_3$ modulo
    $2$, so $K\Gamma_3(X) = -\Gamma_3(X)$. Conversely, if $Kv = -v$, then
    $v(t_1,1,t_3) = -(-1)^{t_1+t_3}v(t_1,0,t_3)$, and the matrix $X$ with
    $X_{0,1-t} = (-1)^t\bigl(v(t,0,0)+v(t,0,1)\bigr)$,
    $X_{1,1-t} = (-1)^t\bigl(v(t,0,0)-v(t,0,1)\bigr)$ satisfies
    $\Gamma_3(X) = 2v$ by (\ref{eq:cluster_source_gamma_three}); hence
    $v = \Gamma_3(X/2)\in\mathcal G_3(A_{\mathrm{src}})$.
\end{proof}

\begin{definition}[Translated cluster stabilizers]\label{def:cluster_chain_stabilizer}
    \lean{MPSTensor.clusterChainStabilizer}
    \leanok
    On a periodic chain of $N$ sites, the stabilizer at site $i$ is
    $K_i = \sigma^z_i\sigma^x_{i+1}\sigma^z_{i+2}$, with indices modulo $N$:
    \begin{align}
        (K_i\psi)(\sigma)
         & = (-1)^{\sigma_i+\sigma_{i+2}}\,
        \psi(\sigma_1,\ldots,\sigma_i,\sigma_{i+1}+1,\sigma_{i+2},\ldots,\sigma_N).
        \notag
    \end{align}
\end{definition}

\begin{theorem}[Cluster stabilizer ground space is the periodic chain ground space]
    \label{thm:cluster_stabilizer_chain_ground_space}
    \lean{MPSTensor.cluster_iInf_eigenspace_eq_chainGroundSpace}
    \leanok
    \uses{def:cluster_tensor, def:chain_ground_space, def:cluster_chain_stabilizer}
    For $N\ge3$, the common $+1$ eigenspace of the translated stabilizers is
    the periodic three-site chain ground space of the cluster tensor:
    \begin{align}
        \bigcap_{i=1}^{N}\{\psi : K_i\psi = \psi\}
         & = \mathcal G_{N,3}(A_{\mathrm{cl}}).
        \notag
    \end{align}
\end{theorem}

\begin{proof}
    \leanok
    \uses{thm:cluster_ground_space_three_eigenspace}
    For a cyclic three-site window at $i$ and an outside assignment $\tau$,
    the restriction $\psi|_{[i,i+2],\tau}$ satisfies
    $K\bigl(\psi|_{[i,i+2],\tau}\bigr) = (K_i\psi)|_{[i,i+2],\tau}$, since
    $K_i$ reads $\sigma_i$, $\sigma_{i+2}$ and flips $\sigma_{i+1}$ inside the
    window. Thus $K_i\psi = \psi$ implies that every restriction is fixed by
    $K$; conversely, if every restriction is fixed by $K$, then evaluating at
    the window extracted from $\sigma$ with outside assignment $\sigma$ itself
    gives $(K_i\psi)(\sigma) = \psi(\sigma)$. By
    Theorem~\ref{thm:cluster_ground_space_three_eigenspace}, a restriction is
    fixed by $K$ exactly when it lies in $\mathcal G_3(A_{\mathrm{cl}})$.
\end{proof}

\begin{theorem}[The cluster state spans the common $+1$ eigenspace of the $ZXZ$ stabilizers]
    \label{thm:cluster_stabilizer_unique_gs}
    \lean{MPSTensor.cluster_iInf_eigenspace_eq_mpvSubmodule, MPSTensor.cluster_stabilizer_unique_gs}
    \leanok
    \uses{def:cluster_tensor, def:mpv_submodule, def:cluster_chain_stabilizer}
    For $N\ge3$,
    \begin{align}
        \bigcap_{i=1}^{N}\{\psi : K_i\psi = \psi\}
         & = \C\,V^{(N)}(A_{\mathrm{cl}}),
        \notag
    \end{align}
    and this common eigenspace is one-dimensional.
\end{theorem}

\begin{proof}
    \leanok
    \uses{thm:cluster_stabilizer_chain_ground_space, thm:cluster_normal, thm:chain_ground_space_eq_mpv_normal, thm:parent_hamiltonian_unique_gs}
    By Theorem~\ref{thm:cluster_stabilizer_chain_ground_space} the common
    eigenspace is $\mathcal G_{N,3}(A_{\mathrm{cl}})$. The cluster tensor is
    injective after blocking $L_0 = 2$ sites
    (Theorem~\ref{thm:cluster_normal}), so
    Theorem~\ref{thm:chain_ground_space_eq_mpv_normal} with $L = 3 \le N$
    gives $\mathcal G_{N,3}(A_{\mathrm{cl}}) = \C\,V^{(N)}(A_{\mathrm{cl}})$,
    and Theorem~\ref{thm:parent_hamiltonian_unique_gs} gives
    $\dim\mathcal G_{N,3}(A_{\mathrm{cl}}) = 1$.
\end{proof}

\begin{remark}[Cluster-state stabilizer convention]\label{rem:cluster_stabilizer_convention}
    The cluster tensor of Definition~\ref{def:cluster_tensor} is the reflected
    convention of the tensor displayed in
    \cite[lines~2364--2369]{Cirac2021Matrix}. By (\ref{eq:cluster_gamma_three})
    and its reflected counterpart, the three-site parent space of either
    convention is the $+1$ eigenspace of $\sigma^z_1\sigma^x_2\sigma^z_3$.
    Since each $K_i$ is a self-adjoint involution, $\tfrac12(\Id-K_i)$ is the
    orthogonal projector onto its $-1$ eigenspace, and
    Theorem~\ref{thm:cluster_stabilizer_unique_gs} says that the cluster state
    is the unique ground state of
    $\sum_i \tfrac12(\Id - \sigma^z_i\sigma^x_{i+1}\sigma^z_{i+2})$, with
    ground energy $0$.
    % The identification of the parent interaction with (1-K)/2 as an
    % operator identity, and the self-adjoint-involution property of K_i,
    % are not formalized; the formal statements are the eigenspace
    % identities above.

    The older review states instead that cluster states are unique ground
    states of $\sum_i \sigma^z_i\sigma^x_{i+1}\sigma^z_{i+2}$ with the sign as
    printed \cite[local TeX lines~374--387]{PerezGarcia2007Matrix}. Both
    statements are correct for their own tensors: on the matrices
    $A_{\mathrm{src}}$ displayed there every stabilizer has eigenvalue $-1$
    (Theorem~\ref{thm:cluster_source_ground_space_three_eigenspace}), while on
    $A_{\mathrm{cl}}$ it has eigenvalue $+1$. The two tensors are related by
    $A_{\mathrm{src}}^t = \sqrt2\,\sigma^x A^t$ for the tensor $A$ of
    \cite[lines~2364--2369]{Cirac2021Matrix}, and since
    $\bra{\pm}\sigma^x = \pm\bra{\pm}$ the states differ by $\sigma^z$ on
    every site, which conjugates each $K_i$ to $-K_i$ and exchanges the two
    Hamiltonians $\sum_i K_i$ and $-\sum_i K_i = \sum_i(\Id-K_i) - N\Id$.
\end{remark}

\begin{theorem}[AKLT two-site parent space]\label{thm:aklt_ground_space_two}
    \lean{MPSTensor.mem_aklt_groundSpace_two_iff}
    \leanok
    \uses{def:aklt_tensor, def:ground_space_parent}
    For the AKLT tensor $A$ of Definition~\ref{def:aklt_tensor}, a two-site
    vector $v$ lies in $\mathcal G_2(A)$ if and only if
    \begin{align}
        v_{11} = v_{22} = 0,
        \quad
        v_{01} + v_{10} = 0,
        \quad
        v_{02} + v_{20} = 0,
        \quad
        2v_{00} + v_{12} + v_{21} = 0.
        \label{eq:aklt_two_site_constraints}
    \end{align}
\end{theorem}

\begin{proof}
    \leanok
    \uses{def:aklt_tensor, def:ground_space_map}
    Write $A^0 = c\,\sigma^z$, $A^1 = b\,\sigma^+$, $A^2 = -b\,\sigma^-$ with
    $c = 1/\sqrt3$, $b = \sqrt2/\sqrt3$, so that $b^2 = 2c^2$. The products are
    $A^0A^0 = c^2\Id$, $A^0A^1 = -A^1A^0 = cb\,\ketbra01$,
    $A^0A^2 = -A^2A^0 = cb\,\ketbra10$, $A^1A^2 = -b^2\ketbra00$,
    $A^2A^1 = -b^2\ketbra11$, and $A^1A^1 = A^2A^2 = 0$. Hence
    $\Gamma_2(X)(a,b) = \tr(A^aA^bX)$ has coordinates
    \begin{align}
        \Gamma_2(X)
         & = \bigl(c^2(X_{00}+X_{11}),\ cbX_{10},\ cbX_{01},\ -cbX_{10},\ 0,\
        -b^2X_{00},\ -cbX_{01},\ -b^2X_{11},\ 0\bigr)
        \label{eq:aklt_gamma_two}
    \end{align}
    in the order $00,01,02,10,11,12,20,21,22$, which satisfies
    (\ref{eq:aklt_two_site_constraints}) because $2c^2 = b^2$. Conversely,
    given $v$ satisfying (\ref{eq:aklt_two_site_constraints}), the matrix
    $X$ with $X_{00} = -c^2v_{12}$, $X_{01} = cb\,v_{02}$, $X_{10} = cb\,v_{01}$,
    $X_{11} = -c^2v_{21}$ satisfies $\Gamma_2(X) = c^2b^2\,v$ by
    (\ref{eq:aklt_gamma_two}), so $v = \Gamma_2(X/(c^2b^2))$.
\end{proof}

\begin{remark}[AKLT two-site parent space as a span]\label{rem:aklt_ground_space_two_span}
    % The spanning set and the dimension count are not formalized; the formal
    % statement is the coordinate description of
    % Theorem~\ref{thm:aklt_ground_space_two}.
    The constraints (\ref{eq:aklt_two_site_constraints}) of
    Theorem~\ref{thm:aklt_ground_space_two} describe $\mathcal G_2(A)$ as the
    four-dimensional span of $\ket{01}-\ket{10}$, $\ket{02}-\ket{20}$,
    $\ket{00}-2\ket{12}$, and $\ket{00}-2\ket{21}$. This span is the spin-0
    and spin-1 subspace of two spin-1 particles.
\end{remark}

\begin{theorem}[AKLT three-site intersection property]
    \label{thm:aklt_ground_space_intersection_two}
    \lean{MPSTensor.aklt_groundSpace_intersection_two, MPSTensor.aklt_mem_groundSpace_three_of_inLeftGround_of_inRightGround}
    \leanok
    \uses{def:aklt_tensor, def:ground_space_parent, def:restrict_last, def:restrict_first}
    For the AKLT tensor $A$, the three-site intersection property holds:
    \begin{align}
        (\mathcal G_{[1,2]}\otimes\mathcal H_3)
        \cap(\mathcal H_1\otimes\mathcal G_{[2,3]})
         & = \mathcal G_{[1,3]},
        \notag
    \end{align}
    where $\mathcal G_{[a,b]}$ denotes the local ground space
    $\mathcal G_{b-a+1}(A)$ placed on sites $a,\ldots,b$. This is the identity
    $\mathcal G_{1,2}\cap\mathcal G_{2,3} = \mathcal G_{1,2,3}$ that
    \cite[line~2095]{Cirac2021Matrix} asks to check by hand; it is not an
    instance of the general intersection property, since the AKLT tensor is
    not injective on one site.
\end{theorem}

\begin{proof}
    \leanok
    \uses{thm:aklt_ground_space_two, lem:ground_space_in_left_ground, lem:ground_space_in_right_ground}
    The inclusion $\supseteq$ is the restriction property of local ground
    spaces (Lemmas~\ref{lem:ground_space_in_left_ground}
    and~\ref{lem:ground_space_in_right_ground}). For $\subseteq$, let $\psi$
    be a three-site vector whose restrictions $\psi(i,\cdot,\cdot)$ and
    $\psi(\cdot,\cdot,k)$ all lie in $\mathcal G_2(A)$, so that each satisfies
    (\ref{eq:aklt_two_site_constraints}). With $c$, $b$ as in the proof of
    Theorem~\ref{thm:aklt_ground_space_two}, the products $A^aA^bA^c$ give
    \begin{align}
        \Gamma_3(X)
         & = c^3(X_{00}-X_{11})\ket{000}
        + c^2b\,X_{10}\bigl(\ket{001}-\ket{010}+\ket{100}\bigr)
        - c^2b\,X_{01}\bigl(\ket{002}-\ket{020}+\ket{200}\bigr)
        \notag                           \\
         & \quad
        + cb^2X_{00}\bigl(\ket{102}-\ket{012}-\ket{120}\bigr)
        + cb^2X_{11}\bigl(\ket{021}-\ket{201}+\ket{210}\bigr)
        - b^3X_{10}\ket{121}
        + b^3X_{01}\ket{212}.
        \label{eq:aklt_gamma_three}
    \end{align}
    Set $X_{00} = c\,\psi_{102}$, $X_{01} = -b\,\psi_{002}$,
    $X_{10} = b\,\psi_{001}$, $X_{11} = c\,\psi_{210}$. Then
    $\Gamma_3(X) = c^2b^2\,\psi$ coordinate by coordinate: the twelve
    coordinates absent from (\ref{eq:aklt_gamma_three}) vanish for $\psi$ by
    the constraints $v_{11} = v_{22} = 0$ on the six restrictions, together
    with $\psi_{101} = -\psi_{110} = 0$ and $\psi_{202} = -\psi_{220} = 0$;
    the sign relations $\psi_{010} = -\psi_{001} = -\psi_{100}$,
    $\psi_{020} = -\psi_{002} = -\psi_{200}$, $\psi_{012} = -\psi_{102} = \psi_{120}$,
    $\psi_{201} = -\psi_{210}$ come from $v_{01}+v_{10} = 0$ and
    $v_{02}+v_{20} = 0$; and the remaining coordinates use
    $2v_{00}+v_{12}+v_{21} = 0$ with $b^2 = 2c^2$:
    $2\psi_{000} = \psi_{102}-\psi_{210}$, $\psi_{021} = \psi_{210}$,
    $\psi_{121} = -2\psi_{001}$, $\psi_{212} = -2\psi_{002}$. Hence
    $\psi = \Gamma_3(X/(c^2b^2)) \in \mathcal G_3(A)$.
\end{proof}

\begin{theorem}[Unique ground state of the two-site AKLT parent Hamiltonian]
    \label{thm:aklt_chain_ground_space_two}
    \lean{MPSTensor.aklt_chainGroundSpace_two_eq_mpvSubmodule, MPSTensor.aklt_parentHamiltonian_two_unique_gs}
    \leanok
    \uses{def:aklt_tensor, def:chain_ground_space, def:mpv_submodule}
    For $N\ge3$, the periodic two-site chain ground space of the AKLT tensor
    is spanned by the AKLT state:
    \begin{align}
        \mathcal G_{N,2}(A_{\mathrm{AKLT}})
         & = \C\,V^{(N)}(A_{\mathrm{AKLT}}).
        \notag
    \end{align}
    Thus the two-site parent Hamiltonian has a unique ground state although
    the AKLT tensor becomes injective only after blocking $L_0 = 2$ sites
    \cite[line~2095]{Cirac2021Matrix}. On two periodic sites the two-site
    chain ground space is $\mathcal G_2(A_{\mathrm{AKLT}})$ itself, so the
    condition $N\ge3$ is necessary.
\end{theorem}

\begin{proof}
    \leanok
    \uses{thm:aklt_ground_space_intersection_two, thm:aklt_normal, thm:chain_ground_space_eq_mpv_normal, lem:mpv_mem_chain_ground_space, thm:mpv_ne_zero_blk_injective}
    Let $\psi\in\mathcal G_{N,2}(A_{\mathrm{AKLT}})$. Every cyclic three-site
    window restricts on either side to a cyclic two-site window, so both
    restrictions of $\psi|_{[i,i+2],\tau}$ lie in
    $\mathcal G_2(A_{\mathrm{AKLT}})$, and
    Theorem~\ref{thm:aklt_ground_space_intersection_two} places
    $\psi|_{[i,i+2],\tau}$ in $\mathcal G_3(A_{\mathrm{AKLT}})$. Hence
    $\mathcal G_{N,2}(A_{\mathrm{AKLT}})\subseteq\mathcal G_{N,3}(A_{\mathrm{AKLT}})$.
    The AKLT tensor is injective after blocking $L_0 = 2$ sites
    (Theorem~\ref{thm:aklt_normal}), so
    Theorem~\ref{thm:chain_ground_space_eq_mpv_normal} with $L = 3\le N$ gives
    $\mathcal G_{N,3}(A_{\mathrm{AKLT}}) = \C\,V^{(N)}(A_{\mathrm{AKLT}})$.
    The reverse inclusion holds because the periodic MPS vector satisfies
    every cyclic two-site constraint
    (Lemma~\ref{lem:mpv_mem_chain_ground_space}). Finally
    $V^{(N)}(A_{\mathrm{AKLT}})\ne0$ by
    Theorem~\ref{thm:mpv_ne_zero_blk_injective} with $L_0 = 2 < N$, so
    $\dim\mathcal G_{N,2}(A_{\mathrm{AKLT}}) = 1$.
\end{proof}

\begin{remark}[AKLT polynomial Hamiltonian and edge modes]\label{rem:aklt_parent_hamiltonian}
    The displayed local spin-chain Hamiltonian of the AKLT state is
    \begin{align}
        H_{\mathrm{AKLT},N}
         & =
        \sum_i
        \left(\vec S_i\cdot\vec S_{i+1}
        +\frac13(\vec S_i\cdot\vec S_{i+1})^2\right),
        \notag
    \end{align}
    up to an additive constant
    \cite[local TeX lines~340--349]{PerezGarcia2007Matrix};
    \cite[lines~2126--2130]{Schuch2011Phases}. The spin-$2$ projector form
    is related to the polynomial Hamiltonian by the spin-addition identity
    \begin{align}
        P^{(2)}_{i,i+1}
         & =
        \frac12
        \left(\vec S_i\cdot\vec S_{i+1}
        +\frac13(\vec S_i\cdot\vec S_{i+1})^2
        +\frac23\Id\right)
        \notag
    \end{align}
    for two spin-$1$ particles. For open boundary conditions and $N\geq2$,
    \begin{align}
        \dim\GS(H_{\mathrm{AKLT},N}^{\mathrm{open}})
         & = 4, \notag \\
        \dim\bigl(\GS(H_{\mathrm{AKLT},N}^{\mathrm{open}})
        \cap\mathcal H_{S=0}
        \bigr)
         & = 1
        \notag
    \end{align}
    \cite[lines~1169--1174]{Cirac2021Matrix}.
    % Neither the polynomial form nor the open-boundary degeneracy is
    % formalized; the periodic two-site uniqueness is
    % thm:aklt_chain_ground_space_two.
\end{remark}
